\documentclass{article}
\usepackage[T1]{fontenc} % makes Times/Courier resolve under XeTeX (tectonic)

% ICML 2026 template (paper/icml2026.sty).  Switch [preprint] -> nothing for
% a blind submission, or [accepted] for camera-ready.
\usepackage{microtype}
\usepackage{graphicx}
\usepackage{booktabs,array}
\usepackage{hyperref}
\newcommand{\theHalgorithm}{\arabic{algorithm}}
\usepackage[preprint]{icml2026}
\usepackage{amsmath,amssymb,mathtools,amsthm}
\usepackage{enumitem}
\usepackage{xcolor}
\usepackage[capitalize,noabbrev]{cleveref}

% ---------------------------------------------------------------------------
% Lean cross-references.  Every \leanref{Name} is checked by CI
% (scripts/check_lean_refs.sh) against lean/ServingQueueTheory/*.lean, and the
% named theorem is audited for `sorry` and non-standard axioms
% (scripts/check_lean.sh).
% ---------------------------------------------------------------------------
\newcommand{\leanref}[1]{\texttt{#1}}
% Footer of a proposition box listing the Lean theorems that prove it.
\newenvironment{leanrefs}{\par\nopagebreak\smallskip
  \noindent\textcolor{propframe!35}{\rule{\linewidth}{0.4pt}}\par\nopagebreak\smallskip
  \begingroup\footnotesize\color{black!55}\raggedright\hyphenpenalty=10000\exhyphenpenalty=10000
  \noindent\mbox{\textsf{\bfseries Verified in Lean:}} \ignorespaces}{\par\endgroup}

% Shaded, breakable boxes for propositions (fonts and margins unchanged, so
% this stays within the ICML style).
\usepackage[most]{tcolorbox}
\definecolor{propframe}{RGB}{40,70,130}
\definecolor{propback}{RGB}{243,246,252}
\tcolorboxenvironment{proposition}{
  enhanced, breakable,
  colback=propback, colframe=propframe,
  boxrule=0pt, leftrule=2pt, arc=0pt, outer arc=0pt,
  left=6pt, right=6pt, top=4pt, bottom=4pt,
  before skip=8pt plus 2pt, after skip=8pt plus 2pt,
}
\newcommand{\lean}[1]{{\small\texttt{#1}}}
\newcommand{\exlean}[1]{\par\nopagebreak{\footnotesize\color{black!55}\hyphenpenalty=10000\exhyphenpenalty=10000\raggedleft #1\par}}

\theoremstyle{plain}
\newtheorem{theorem}{Theorem}[section]
\newtheoremstyle{propstyle}{0pt}{0pt}{\normalfont}{}{\bfseries\color{propframe}}{.}{ }{\thmname{#1}\thmnumber{ #2}\thmnote{ (#3)}}
\theoremstyle{propstyle}
\newtheorem{proposition}[theorem]{Proposition}
\theoremstyle{plain}
\newtheorem{lemma}[theorem]{Lemma}
\newtheorem{corollary}[theorem]{Corollary}
\theoremstyle{definition}
\newtheorem{definition}[theorem]{Definition}
\newtheorem{example}[theorem]{Example}
\theoremstyle{remark}
\newtheorem{remark}[theorem]{Remark}

\newcommand{\E}{\mathbb{E}}
\newcommand{\Prob}{\mathbb{P}}
\DeclareMathOperator*{\argmin}{arg\,min}

% icml2026.sty rejects running heads whose box height exceeds 6.25pt; with
% tectonic/XeTeX metrics every \small\bf line is slightly taller, so we cap the
% reported height explicitly.
\icmltitlerunning{\raisebox{0pt}[6pt][0pt]{Decision-Faithful Queueing Models for Agentic LLM Serving}}

\begin{document}

\twocolumn[
  \icmltitle{Decision-Faithful Queueing Models for Agentic LLM Serving}

  \begin{icmlauthorlist}
    \icmlauthor{Jinhwan Suk}{rbln}
  \end{icmlauthorlist}

  \icmlaffiliation{rbln}{Rebellions Inc., Seongnam, Republic of Korea}
  \icmlcorrespondingauthor{Jinhwan Suk}{jinhwan.suk@rebellions.ai}

  \icmlkeywords{LLM serving, queueing theory, agentic workloads, KV cache,
    prefill/decode disaggregation, formal verification, Lean}

  \vskip 0.3in
]

\printAffiliationsAndNotice{Working draft v0.1. Analytical results are
machine-checked; empirical results are pending.}

\begin{abstract}
Benchmarks of LLM serving systems report kernel latency and tokens per
second, but service-level latency is dominated by queueing once utilisation
approaches saturation. We argue that a layered analytical model --- queueing
theory for congestion, caching theory for KV state locality, queueing
networks for multi-turn agent programs, and stochastic control for
routing/eviction decisions --- turns measurement into prediction and
lets design decisions be derived from one cost function instead of a
collection of heuristics. We state seven propositions that this framework
yields and \emph{machine-check every one of them in Lean~4/Mathlib};
the proofs are compiled, audited for unfinished goals and non-standard
axioms, and cross-referenced with this document in continuous integration.
Using the propositions we re-examine three claims made about agentic
serving: that KV offloading is ineffective, that prefill/decode (PD)
disaggregation is ineffective, and that shortest-context-first eviction is
optimal. We show the first two hold only in identifiable operating
regimes (we give the exact win condition for PD), and that the third is
false in general (a three-program counterexample; the problem is a
covering knapsack). Finally we propose a validation methodology in which
the target is \emph{decision faithfulness} --- agreement on which policy
is better --- rather than cycle-accurate reproduction, and an empirical
plan for a rack-scale NPU serving testbed.
\end{abstract}

% ===========================================================================
\section{Introduction}
% ===========================================================================

Response time of a serving system decomposes as
\[
  T = W_q + S ,
\]
queueing delay plus service time. Hardware benchmarks measure~$S$; users
experience~$T$. Near saturation $W_q$ dominates, and it does so
\emph{non-linearly}: in the simplest Markovian model a 10\,\% increase in
arrival rate (from $\rho = 0.90$ to $\rho = 0.99$) multiplies response time
by ten (Proposition~\ref{prop:mm1}). This single fact reframes most serving
questions: ``how many tokens/s does this NPU deliver?'' becomes ``at what
utilisation does p99~TTFT cross the SLA, and which control knob moves that
point?''

Agentic workloads sharpen the problem. In production coding-agent traces the median
session has 43 LLM turns, a median input of 142K tokens, a median output of
444 tokens and a prefix-cache hit rate above 96\,\% \citep{vllm-agentx}:
each turn appends a few hundred tokens to a context of $10^5$ tokens,
interleaved with tool calls of unpredictable duration. The unit of work that matters is the \emph{program}
(a multi-turn session with persistent KV state), not the request
\citep{thunderagent}. Whether a program's KV cache is kept, offloaded, or
recomputed, and on which replica the next turn runs, determines the
service-time distribution that feeds the queue.

This paper makes three contributions.

\begin{enumerate}[leftmargin=*]
\item \textbf{A layered model} (\S\ref{sec:model}) in which queueing,
caching, program state and control are separate layers with explicit
interfaces: the eviction policy sets the hit rate; the hit rate sets the
service-time distribution; the distribution sets utilisation and delay;
delay enters the routing/eviction cost.
\item \textbf{Seven machine-checked propositions}
(\S\ref{sec:props}). Each is stated informally here and proved formally in
Lean~4 against Mathlib. We do this because the propositions were
\emph{drafted with the help of a large language model}, and LLM-produced
``proofs'' are not trustworthy without independent checking. The CI
pipeline compiles the proofs, rejects any \lean{sorry}, prints the axioms
each theorem depends on and fails if anything beyond \lean{propext},
\lean{Classical.choice}, \lean{Quot.sound} appears, and checks that every
\lean{\textbackslash leanref} in this document names an existing, audited
theorem.
\item \textbf{A re-examination of three published claims}
(\S\ref{sec:thunder}) and a \textbf{validation methodology}
(\S\ref{sec:validation}) whose success criterion is decision faithfulness.
\end{enumerate}

What this paper does \emph{not} do: it does not yet report measurements.
\S\ref{sec:plan} lays out the empirical programme; the analytical results
are meant to tell us which measurements are worth taking.

% ===========================================================================
\section{A Layered Model of Agentic Serving}\label{sec:model}
% ===========================================================================

\subsection{Layer 1: queueing}

Requests arrive at rate~$\lambda$ to a service centre (one replica, one
device pool) with service time~$S$. Utilisation is $\rho=\lambda\E[S]$
(single server). For M/M/1 the mean response time is
$W = 1/(\mu-\lambda)$ with $\mu = 1/\E[S]$; for M/G/1 the
Pollaczek--Khinchine (PK) formula \citep{kleinrock1975} gives
\begin{equation}\label{eq:pk}
  \E[W_q] = \frac{\lambda \E[S^2]}{2(1-\rho)} .
\end{equation}
The second moment $\E[S^2] = \mathrm{Var}[S] + \E[S]^2$ makes
\emph{service-time variance} a first-class quantity
(Proposition~\ref{prop:pk}). Where that variance comes from is workload
dependent, and we do not assume agentic workloads are intrinsically
high-variance. In a steady-state agentic deployment with a prefix hit rate
above 96\,\% \citep{vllm-agentx}, a turn is an append-prefill of a few
hundred tokens plus a decode of a few hundred tokens; if every turn hit,
$\E[S^2]$ would be close to $\E[S]^2$. The variance enters through the
\emph{hit/miss mixture}: a miss re-prefills a $10^5$-token context and
costs one to two orders of magnitude more than a hit, and a two-point
mixture with a 4\,\% miss probability and a $100\times$ miss penalty has
squared coefficient of variation $\mathrm{CV}^2>15$
(Example~\ref{ex:cv2}). The M/G/1 delay is $(1+\mathrm{CV}^2)/2$ times the
M/M/1 delay at the same load, so eviction policy is what sets the variance
term. A second, fleet-level source is mixing request classes (short chat,
long document, agent) in one FIFO queue; that is the classical root of
class separation and SRPT-style scheduling \citep{schrage1968}. Which
source dominates in a given deployment is an empirical question
(\S\ref{sec:plan}).

\subsection{Layer 2: caching}

A KV hit changes the service time:
\begin{gather*}
  S = \begin{cases} S^{\mathrm{hit}} & \text{w.p. } p \\
                    S^{\mathrm{miss}} & \text{w.p. } 1-p, \end{cases}\\
  \E[S] = p\,S^{\mathrm{hit}} + (1-p)\,S^{\mathrm{miss}} .
\end{gather*}
The eviction/admission policy determines~$p$, hence $\E[S]$, $\E[S^2]$
and~$\rho$. Because $W_q$ is convex in~$\rho$ and unbounded as
$\rho\to1$, cache reuse is not merely a compute saving; it is a
\emph{queue-emptying} mechanism (Proposition~\ref{prop:cache}). The
right objective for an eviction policy is therefore not hit ratio but
$\min \E[\mathrm{TTFT}]$ or $\min \Prob(\mathrm{TTFT} > \mathrm{SLA})$,
and a natural per-byte utility is
\begin{equation}\label{eq:utility}
  U_i = \frac{\Prob(\text{reuse}_i)\cdot \Delta S_i}{\mathrm{KVBytes}_i}\cdot
        \frac{1}{1-\rho},
\end{equation}
where the last factor makes hot state more valuable as the system
approaches saturation.

\subsection{Layer 3: programs as jobs in a queueing network}

An agent program cycles through
$\text{prefill}\to\text{decode}\to\text{tool}\to\text{prefill}\to\cdots$.
Tool execution plays the role of \emph{think time} in an interactive
closed network \citep{lazowska1984}. A program in the tool phase is not
``finished''; it is suspended future work with resident state. The state
$z_p(t)\in\{\text{reasoning},\text{acting},\text{waiting}\}$ and the
transition probabilities $\Prob(\text{next phase}\mid z)$ are observable,
and they are exactly what a request-level scheduler discards.

\subsection{Layer 4: stochastic control}

Let $X_t = (Q_t, \mathrm{KV}_t, P_t, R_t)$ collect queue states, KV
placement, program states, and resources; let actions be
$A_t \in \{\text{route},\text{evict},\text{offload},\text{migrate},
\text{admit},\text{prefetch}\}$. The serving problem is
\begin{multline}\label{eq:mdp}
  \min_\pi \E_\pi\Bigl[\sum_t
     \alpha\,\text{latency}_t + \beta\,\text{recompute}_t \\
   + \gamma\,\text{migration}_t + \delta\,\text{SLA-violation}_t \Bigr].
\end{multline}
A myopic router minimises $W_j + S_{ij}$; a KV-aware router adds a
locality term; a program-aware router adds migration cost and
$\E[V(X_{t+1})\mid p,j]$. Session affinity, KV-aware routing and
PPD-style append-prefill routing are all one-step approximations of
Eq.~\eqref{eq:mdp} (Proposition~\ref{prop:routing}).

% ===========================================================================
\section{Machine-Checked Propositions}\label{sec:props}
% ===========================================================================

Each shaded box below is a statement that has been formally proved in the
accompanying Lean~4 development;\footnote{\texttt{lean/ServingQueueTheory/}
in the artifact. CI compiles the proofs, rejects \lean{sorry}, prints the
axioms of every theorem named in this paper and fails on anything beyond
\lean{propext}, \lean{Classical.choice}, \lean{Quot.sound}, and checks that
every theorem name cited here exists. CI does not check that the Lean
statement faithfully formalises the prose; the statements were kept close
to the prose to make that human review easy.} its footer names the Lean
theorems. Statements in the boxes are informal paraphrases; the Lean
statement is authoritative. Human-readable proofs are in
Appendix~\ref{app:proofs}. Numerical instances are
separated into \emph{Examples}; these too are checked (by \lean{norm\_num}
or \lean{decide}).

\begin{proposition}[M/M/1 latency blow-up]\label{prop:mm1}
Let $\mu>0$ and $W(\lambda)=1/(\mu-\lambda)$ on $0\le\lambda<\mu$, with
$\rho=\lambda/\mu$. Then
\begin{enumerate}[nosep,leftmargin=1.6em,label=(\roman*)]
\item $W(\lambda) = \dfrac{1/\mu}{1-\rho}$;
\item $W$ is strictly increasing on $[0,\mu)$;
\item $W$ is unbounded: for every $M$ there is a stable load $\lambda<\mu$
      with $W(\lambda)>M$.
\end{enumerate}
\begin{leanrefs}\leanref{mm1Wait\_eq\_rho\_form}, \leanref{mm1Wait\_strictMono}, \leanref{mm1Wait\_unbounded}\end{leanrefs}
\end{proposition}

\begin{example}\label{ex:mm1}
With $\mu=10$\,req/s, arrival rates $9$, $9.5$, $9.9$\,req/s give
$W=1$, $2$, $10$\,s: a 10\,\% increase in load multiplies latency by ten.
\exlean{\leanref{mm1Wait\_example\_ratio}}
\end{example}

\begin{proposition}[Variance drives queueing delay]\label{prop:pk}
Let $S$ have a finite discrete distribution and let
$\E[W_q]=\lambda\E[S^2]/(2(1-\rho))$ be the PK delay~\eqref{eq:pk}.
\begin{enumerate}[nosep,leftmargin=1.6em,label=(\roman*)]
\item $\E[S^2]=\mathrm{Var}[S]+\E[S]^2$;
\item for fixed $\lambda>0$ and $\rho<1$, $\E[W_q]$ is strictly increasing
      in $\E[S^2]$;
\item if two workloads have equal mean and $\mathrm{Var}[S_A]<\mathrm{Var}[S_B]$,
      then $\E[W_q^A]<\E[W_q^B]$ at every stable load.
\end{enumerate}
\begin{leanrefs}\leanref{secondMoment\_eq\_variance\_add\_sq}, \leanref{pkWait\_strictMono\_secondMoment}, \leanref{pkWait\_lt\_of\_variance\_lt}\end{leanrefs}
\end{proposition}

\begin{example}\label{ex:pk}
$A=\{1000,1000,1000,1000\}$ and $B=\{100,100,100,3700\}$ have equal mean
$1000$ but $\E[S_A^2]=1.00\times10^6$, $\E[S_B^2]=3.43\times10^6$; hence
$\E[W_q^B]=3.43\,\E[W_q^A]$ for every $\lambda$ and $\rho<1$.
\exlean{\leanref{workloadB\_wait\_ratio}}
\end{example}

\begin{proposition}[Cache reuse empties the queue]\label{prop:cache}
Let $0\le S^{\mathrm{hit}}\le S^{\mathrm{miss}}$ and let a request hit with
probability~$p$, so $\E[S]=pS^{\mathrm{hit}}+(1-p)S^{\mathrm{miss}}$ and
$\rho=\lambda\E[S]$. Then
\begin{enumerate}[nosep,leftmargin=1.6em,label=(\roman*)]
\item $\E[S]$, $\E[S^2]$ and $\rho$ are antitone in $p$;
\item the PK delay of the hit/miss mixture is antitone in $p\in[0,1]$
      whenever the system is stable at the lower hit rate.
\end{enumerate}
\begin{leanrefs}\leanref{meanService\_antitone}, \leanref{utilization\_antitone}, \leanref{secondMomentService\_antitone}, \leanref{pkWait\_mixture\_antitone}\end{leanrefs}
\end{proposition}

\begin{example}\label{ex:cache}
$S^{\mathrm{hit}}=50$\,ms, $S^{\mathrm{miss}}=500$\,ms, $\lambda=1.8$\,req/s.
At $p=0$: $\rho=0.90$. At $p=0.8$: $\E[S]=140$\,ms and $\rho=0.252$. By
Proposition~\ref{prop:mm1}, the second point is far from the blow-up
region; the first is inside it.
\exlean{\leanref{utilization\_example\_hit80}}
\end{example}

\begin{example}[Where the variance comes from]\label{ex:cv2}
For any mean and load, $\E[W_q]_{M/G/1}/\E[W_q]_{M/M/1}=(1+\mathrm{CV}^2)/2$
with $\mathrm{CV}^2=\E[S^2]/\E[S]^2-1$. With a 96\,\% hit rate,
$S^{\mathrm{hit}}=50$\,ms and $S^{\mathrm{miss}}=5$\,s (full re-prefill of a
long context), $\mathrm{CV}^2>15$: more than $8\times$ the delay of
exponential service. With the same hit rate but a cheap miss
($S^{\mathrm{miss}}=100$\,ms), $\mathrm{CV}^2<0.05$. The variance is a
property of the miss penalty and the eviction policy, not of the workload
label.
\exlean{\leanref{pkWait\_ratio\_to\_exponential}, \leanref{mixtureCV2\_agentic\_example}, \leanref{mixtureCV2\_cheap\_miss\_example}}
\end{example}

\begin{proposition}[Option value of a mechanism]\label{prop:option}
For finite action sets $A_0\subseteq A_1$ and any cost $J$,
\[
  \min_{a\in A_1}J(a)\;\le\;\min_{a\in A_0}J(a).
\]
In particular, with $A_0=\{\text{keep},\text{recompute}\}$ and
$A_1=A_0\cup\{\text{offload}\}$: \emph{enabling} offloading never increases
the optimal cost, while there exist costs $J$ for which
\emph{always} offloading is strictly worse than the $A_0$ optimum.
\begin{leanrefs}\leanref{optimal\_cost\_antitone\_in\_actions}, \leanref{enabling\_offload\_never\_hurts}, \leanref{always\_offload\_can\_be\_worse}\end{leanrefs}
\end{proposition}

\begin{proposition}[PD disaggregation: no-gain theorem and exact win condition]\label{prop:pd}
Let $N$ identical devices serve requests with prefill work $s_P>0$ and
decode work $s_D>0$ (device-seconds). Ideal aggregation has capacity
$\mu_A=N/(s_P+s_D)$; a static split $N_P+N_D=N$ has capacity
$\mu_{PD}=\min(N_P/s_P,\,N_D/s_D)$.
\begin{enumerate}[nosep,leftmargin=1.6em,label=(\roman*)]
\item $\mu_{PD}\le\mu_A$ for every split.
\item Equality holds at the rate-matched split $N_P=Ns_P/(s_P+s_D)$,
      which is therefore optimal.
\item Let aggregation pay an interference overhead~$I$, let dedicated pools
      realise specialisation gains $g_P,g_D$, and let PD be further capped
      by KV-transfer capacity $B_{\mathrm{net}}/\E[K]$ and memory caps
      $\mu_{P,\mathrm{mem}},\mu_{D,\mathrm{mem}}$. Then
      \begin{multline*}
        \mu_{PD} > \mu_A \iff
        \frac{N}{s_P+s_D+I} < \\
        \min\Bigl[\tfrac{N}{s_P/g_P+s_D/g_D},\;
                  \tfrac{B_{\mathrm{net}}}{\E[K]},\;
                  \mu_{P,\mathrm{mem}},\;\mu_{D,\mathrm{mem}}\Bigr],
      \end{multline*}
      i.e.\ PD wins iff \emph{each} of its four bottlenecks exceeds the
      aggregated capacity. With $g_P=g_D=1$, $I=0$ the compute term
      collapses to $\mu_A$, recovering~(i).
\end{enumerate}
\begin{leanrefs}\leanref{pd\_le\_agg}, \leanref{pd\_eq\_agg\_at\_rate\_match}, \leanref{rate\_match\_optimal}, \leanref{pd\_beats\_agg\_iff}, \leanref{pdCompute\_eq\_agg\_of\_no\_gain}\end{leanrefs}
\end{proposition}

\begin{example}\label{ex:pd}
$N=32$, $s_P=s_D=1$, $I=0.5$, $g_P=2$, $g_D=1$, ample memory.
PD wins when $B_{\mathrm{net}}/\E[K]=1000$ and loses when
$B_{\mathrm{net}}/\E[K]=10$: the same compute configuration lands on
either side of the condition depending on the KV-transfer budget.
\exlean{\leanref{pd\_wins\_example}, \leanref{pd\_loses\_example}}
\end{example}

\begin{proposition}[Shortest-context-first eviction is not optimal]\label{prop:evict}
Consider suspended programs with context lengths $c_i$, recompute cost
$R(c)=c^2$ and memory target $\Delta C$:
$\min_{S}\sum_{i\in S}c_i^2$ s.t.\ $\sum_{i\in S}c_i\ge\Delta C$.
\begin{enumerate}[nosep,leftmargin=1.6em,label=(\roman*)]
\item The statement ``for every sorted instance and every feasible $S$,
      the shortest-first choice costs no more than $S$'' is \emph{false}:
      for $c=\{4,5,6\}$, $\Delta C=6$, shortest-first evicts $\{4,5\}$
      (cost~41) while $\{6\}$ is feasible with cost~36.
\item With resume probabilities $p_i$, the expected cost is
      $p_ic_i^2=(p_ic_i)\cdot c_i$, and shortest-first fails even for a
      single eviction: $(c,p)=(2,1)$ costs $4$, $(10,0.01)$ costs $1$.
\end{enumerate}
\begin{leanrefs}\leanref{shortestFirst\_not\_optimal}, \leanref{shortestFirst\_optimality\_claim\_false}, \leanref{shortestFirst\_wrong\_with\_resume\_prob}\end{leanrefs}
\end{proposition}

\begin{proposition}[Program-aware routing and the limit of affinity]\label{prop:routing}
Let node~1 hold a program's state and node~2 be a candidate, with queueing
delays $W_j$, service times $S_j$, migration cost $M_2$ and expected future
costs $F_j$.
\begin{enumerate}[nosep,leftmargin=1.6em,label=(\roman*)]
\item The lookahead rule $W_j+S_j+M_j+F_j$ prefers node~1 iff
      $(W_1+S_1)-(W_2+S_2) < M_2+(F_2-F_1)$.
\item For \emph{any} finite $M,F\ge0$ there is a stable load $\lambda<\mu$
      at which the affinity node's M/M/1 wait exceeds $1/\mu+M+F$; hence
      strict session affinity is not an optimal policy.
\item (Append-prefill) Decode-local processing of turn $t\ge2$ beats the
      prefill-pool path iff $S_D+I < W_P+S_P+T_{\mathrm{KV}}$.
\end{enumerate}
\begin{leanrefs}\leanref{lookahead\_prefers\_iff}, \leanref{affinity\_not\_always\_optimal}, \leanref{append\_prefill\_rule}\end{leanrefs}
\end{proposition}

\begin{example}\label{ex:routing}
Node~1 queue 300\,ms with the KV resident; node~2 queue 100\,ms but a
500\,ms KV transfer. Myopic routing picks node~2; with equal futures the
lookahead rule picks node~1. For a turn with a 100K cached prefix and 500
new tokens, decode-local append-prefill (50+20+30\,ms) beats the
prefill-pool path (10+20+400+50\,ms).
\exlean{\leanref{affinity\_example}, \leanref{append\_prefill\_example}}
\end{example}

\begin{remark}
Propositions~\ref{prop:mm1}--\ref{prop:cache} are textbook; we include
them because the downstream arguments depend on their exact form (e.g.\
unboundedness in Prop.~\ref{prop:mm1}(iii) is what Prop.~\ref{prop:routing}(ii)
uses to bound affinity). Propositions~\ref{prop:option}--\ref{prop:routing}
are the ones that bear on published claims.
\end{remark}

% ===========================================================================

% ===========================================================================
\section{Re-examining Three Claims about Agentic Serving}\label{sec:thunder}
% ===========================================================================

ThunderAgent \citep{thunderagent} introduces program-aware scheduling.
Its Appendix~A.2 reports (Fig.~7) that under its workload
(a)~vLLM$+$LMCache offloading collapses at high concurrency --- 111, 54,
62, 40 step/min with 88\,\%, 2\,\%, 1\,\%, 1\,\% hit rate at 24, 48, 72, 96
parallel programs, attributed to PCIe bandwidth --- and (b)~static PD
splits 1:3, 2:2, 3:1 reach 87.1, 160.1, 183.2 step/min against 214.9 for
aggregated serving, attributed to a smaller effective HBM pool for prefill
(``PD disaggregation exacerbates thrashing''). Lemma~4.1 models recompute
cost as quadratic in context length, Definition~4.1 poses
$\min_S\sum_{i\in S}c_i^2$ s.t.\ $\sum_{i\in S}c_i\ge\Delta C$, and
Appendix~F.3 states a theorem that selecting the programs with the
smallest~$c_i$ is optimal. \emph{We have not re-run these experiments; the
discussion below is about what the reported numbers and stated results can
and cannot support.}

\paragraph{Offloading.} The observation is consistent with our model: an
offloading tier is a second service centre with its own capacity
$B_{\mathrm{tier}}/\E[K]$, and once
$\lambda\,\Prob(\text{offload})\,\E[K] > B_{\mathrm{tier}}$ it saturates
and every miss pays the transfer. But by Proposition~\ref{prop:option},
this is evidence that \emph{always-offload} is a bad policy at that
operating point, not that \emph{enabling} offloading is harmful. The
optimal controller chooses $a_i^\ast = \argmin\{J_{\mathrm{keep}},
J_{\mathrm{offload}}, J_{\mathrm{recompute}}\}$ per program and can never do
worse than never offloading. Version~3 of the paper reaches the
same conclusion: Appendix~A.4 (added in v3) states that ThunderAgent ``is
orthogonal to KV-cache offloading'', demonstrates it on SGLang HiCache, and
reframes the finding as ``offloading alone is insufficient at high
concurrency''. The research question is therefore
\emph{program-aware selective offloading}, and its natural objective is
Eq.~\eqref{eq:utility} evaluated per program.

\paragraph{PD disaggregation.} Proposition~\ref{prop:pd}(i)--(ii) says the
reported result is \emph{exactly what the model predicts} under the
paper's conditions: same hardware, no phase specialisation
($g_P=g_D=1$), and a memory-bound prefill pool (a smaller HBM pool
thrashes sooner, i.e.\ $\mu_{P,\mathrm{mem}}$ binds). Nothing about
agentic workloads is needed for this; the bound is a rate-matching
identity. Proposition~\ref{prop:pd}(iii) gives the precise condition
under which the conclusion flips. The regime in which PD is deployed
in practice \citep{vllm-agentx,mooncake-vllm} --- long-context prefill
parallelism ($g_P>1$), wide-EP decode ($g_D>1$), measurable
prefill/decode interference ($I>0$), fast interconnect
($B_{\mathrm{net}}/\E[K]$ not binding), and a shared KV tier that removes
$\mu_{P,\mathrm{mem}}$ as a constraint --- is exactly the regime the
condition identifies. The correct statement is regime-dependent:
``static PD without specialisation and with a memory-bound prefill pool
cannot beat aggregation'' (a theorem) rather than ``PD does not help
agentic serving'' (a generalisation the data does not support).

Moreover, for turn $t\ge2$ the request shape is ``$10^5$ cached tokens
$+\,500$ new tokens''. Routing it through the prefill pool and shipping
the whole KV is dominated whenever
$S_D + I < W_P + S_P + T_{\mathrm{KV}}$ (Prop.~\ref{prop:routing}), which
is the PPD append-prefill rule \citep{ppd}. Pure PD is thus the wrong
abstraction for multi-turn programs; \emph{dynamic} PD with per-turn
routing is what the cost model yields.

\paragraph{Shortest-context-first eviction.} Proposition~\ref{prop:evict}
refutes the Appendix~F.3 theorem as stated. Its proof is an exchange
argument on $(c_{\mathrm{short}}+r)^2 > c_{\mathrm{short}}^2 + r^2$ that
replaces a long program by a short one ``and theoretically the residue
$r$''; but the decision variable is a \emph{subset of programs}, and the
residue is not a program that can be selected. The argument is valid for
the divisible relaxation, not for Definition~4.1. The problem is a covering (min-)knapsack
\citep{csirik1991,carnes2008}, NP-hard in general but with an FPTAS, and shortest-first is a heuristic
whose approximation ratio is unbounded (take $c=\{k,\dots,k, K\}$ with
$\Delta C = K$ and $k\ll K$). With resume probabilities the fractional
relaxation is solved by evicting in increasing order of the density
$p_i c_i$, not $c_i$. The paper's own ablation (Appendix~G.3, v3) shows shortest-first
beating longest-first and random by about $2\times$ in throughput, which is
consistent with shortest-first being a good heuristic; it does not test
against the covering-knapsack optimum. Whether the gap matters in practice
is an empirical question we include in \S\ref{sec:plan}.

% ===========================================================================
\section{Validation Methodology: Decision Faithfulness}\label{sec:validation}
% ===========================================================================

An analytical model of a serving system will not match the real system
cycle for cycle. The question is \emph{which discrepancies matter for the
decisions the model is used to make}. We propose the following protocol.

\subsection{Three levels of validity}
\begin{enumerate}[leftmargin=*]
\item \textbf{Structural validity.} The mechanisms that drive decisions
are represented: queue state, KV locality and transfer cost, memory
pressure, program phase transitions, service-time dependence on cached
tokens.
\item \textbf{Calibration.} Service and transfer models are fitted from
profiling, not assumed:
$S_{\mathrm{prefill}}(L,K,B)$, $S_{\mathrm{decode}}(B,\mathrm{KV})$,
$T_{\mathrm{transfer}}(\mathrm{KVBytes})$,
$\Prob(\text{next phase}\mid z)$.
\item \textbf{Held-out validation.} Compare on workloads not used for
calibration, along three axes: \emph{metric accuracy}
(MAPE of TTFT, TPOT, p99, throughput), \emph{trend accuracy} (sign of the
response to load, hit rate, tool time), and \emph{decision accuracy}
(does the model pick the same routing/eviction/split as the real optimum;
does it rank policies in the same order).
\end{enumerate}

\subsection{Success criteria}
Decision accuracy is primary. A model that predicts a TTFT reduction of
$430\to370$\,ms when the real system shows $500\to420$\,ms is 14\,\% off
in level but right about the decision, and that is what a scheduler
consumes. We therefore report policy-ranking agreement (Kendall~$\tau$
over policy sets) and optimal-decision agreement (fraction of states where
$\argmin$ coincides) alongside MAPE.

\subsection{Error decomposition and ablation}
Report
$\mathrm{Error} = E_{\text{service}} + E_{\text{queue interaction}} +
E_{\text{network}} + E_{\text{program prediction}}$ by substituting
measured quantities for modelled ones one layer at a time. Ablate each
cost term (queue, KV locality, future program state, memory pressure) and
show the decision accuracy lost, so that a claim such as ``the
program-aware term is necessary'' is backed by a measured delta.

\subsection{Out-of-distribution checks}
Vary tool time ($2\,\mathrm{s}\to10\,\mathrm{s}$), context skew, KV-reuse
rate, and program-population balance beyond the calibration range and
report whether policy rankings persist.

\subsection{Scope statement}
The model is not intended to reproduce collective-communication
contention, kernel-launch jitter, allocator fragmentation or network
head-of-line blocking; in regimes where those dominate, error grows and
we say so. A model whose failure regions are known is more useful than
one that claims none.

% ===========================================================================
\section{Empirical Research Plan}\label{sec:plan}
% ===========================================================================

The plan targets a rack-scale NPU deployment; all steps are
hardware-agnostic.

\begin{enumerate}[leftmargin=*]
\item \textbf{Service-time models.} Profile
$S_{\mathrm{prefill}}(L,K,B)$ on a grid of prompt length~$L$, cached
prefix~$K$ and batch~$B$; $S_{\mathrm{decode}}(B,\mathrm{KV})$;
$T_{\mathrm{transfer}}$ vs.\ KV bytes over the actual interconnect.
Fit and report residuals.
\item \textbf{Variance experiment (Prop.~\ref{prop:pk}).} First
\emph{measure} the per-turn service-time $\mathrm{CV}^2$ on replayed agent
traces and decompose it into hit/miss mixture, output-length spread, and
context-length spread; we expect the first to dominate and the others to
be modest, but this is the claim to test. Then vary the hit rate (by
controlling eviction) and check that $W_q$ tracks $(1+\mathrm{CV}^2)/2$.
For fleet-level mixing, compare a shared FIFO against class-separated
queues on chat$+$agent traces of equal mean.
\item \textbf{Cache-reuse experiment (Prop.~\ref{prop:cache}).} Sweep hit
rate by controlling prefix sharing; measure $\rho$ and p99~TTFT; compare
the per-program utility~\eqref{eq:utility} against LRU and hit-ratio
maximisation as eviction objectives.
\item \textbf{Selective offloading (Prop.~\ref{prop:option}).}
Always-offload vs.\ never-offload vs.\ per-program $\argmin$ over
$\{\text{keep},\text{offload},\text{recompute}\}$; sweep concurrency past the
tier's bandwidth saturation and confirm the enabled optimum never falls
below never-offload.
\item \textbf{PD win condition (Prop.~\ref{prop:pd}).} Measure $s_P, s_D,
I, g_P, g_D, B_{\mathrm{net}}/\E[K]$ independently, evaluate the
condition, then run aggregated vs.\ rate-matched PD and check that the
inequality predicts the winner. Repeat with a shared KV tier and with
per-turn append-prefill routing.
\item \textbf{Eviction gap (Prop.~\ref{prop:evict}).} On replayed agent
traces, compare shortest-first against a covering-knapsack solver (exact
for small~$n$, greedy-by-density otherwise) and against the
$p_ic_i$-density rule using empirically estimated resume probabilities.
Report recompute cost and TTFT.
\item \textbf{Routing (Prop.~\ref{prop:routing}).} Myopic, KV-aware,
program-aware (one-step lookahead), and strict affinity; measure the load
at which affinity's advantage inverts and compare with the prediction.
\item \textbf{Decision-faithfulness scorecard} per
\S\ref{sec:validation}, with error decomposition and ablations.
\end{enumerate}

% ===========================================================================
\section{Related Work}
% ===========================================================================

Queueing foundations: \citet{kleinrock1975}; PK formula and M/G/1
\citep{pollaczek1930,khinchine1932}; SRPT optimality
\citep{schrage1968}; interactive closed networks
\citep{lazowska1984}. Covering knapsack: \citet{csirik1991,carnes2008}.
Agentic serving: ThunderAgent \citep{thunderagent}; PPD append-prefill
\citep{ppd}; production traces and PD rate matching
\citep{vllm-agentx}; distributed KV pools \citep{mooncake-vllm}.

% ===========================================================================
\section{Limitations}
% ===========================================================================

The propositions are about the model, not the world; their empirical
force depends on the calibration and validation programme of
\S\S\ref{sec:validation}--\ref{sec:plan}, which has not yet been carried
out. Propositions~\ref{prop:mm1}--\ref{prop:cache} assume Poisson
arrivals and a single server; agentic arrivals are bursty and correlated
(turns of the same program are not independent), which the PK formula
underestimates. We have not measured the service-time variance of a real
agentic deployment; \S\ref{sec:model} argues where it should come from, and
\S\ref{sec:plan} makes measuring it the first experiment. The PD model treats capacity as a scalar bottleneck and
ignores batching dynamics inside each pool. The ThunderAgent statements we
discuss were checked against arXiv:2602.13692v3; we have not reproduced its
experiments.

\section*{Impact Statement}
This work provides analytical tools for capacity planning and scheduling in
LLM serving systems and a workflow for machine-checking LLM-assisted
mathematical claims. Better utilisation of serving hardware lowers the energy
and cost of deploying language models; we see no direct negative societal
consequences beyond those generally associated with more efficient LLM
deployment.

\bibliographystyle{icml2026}
\bibliography{refs}

\appendix
% ===========================================================================
\section{Proofs}\label{app:proofs}
% ===========================================================================

\sloppy
Each proof below is a transcription of the corresponding Lean proof;
the Lean name is given at the end of each item.

\begin{proof}[Proof of Proposition~\ref{prop:mm1}]
(i) $\frac{1}{\mu-\lambda}=\frac{1/\mu}{1-\lambda/\mu}$ since $\mu\neq0$
(\lean{mm1Wait\_eq\_rho\_form}).
(ii) For $0\le a<b<\mu$ we have $0<\mu-b<\mu-a$, hence
$\frac1{\mu-a}<\frac1{\mu-b}$ (\lean{mm1Wait\_strictMono}).
(iii) Given $M$, let $\varepsilon=\min\bigl(\mu,\tfrac{1}{|M|+1}\bigr)>0$
and $\lambda=\mu-\varepsilon\in[0,\mu)$. Then
$W(\lambda)=1/\varepsilon\ge|M|+1>M$ (\lean{mm1Wait\_unbounded}).
Example~\ref{ex:mm1} is arithmetic (\lean{mm1Wait\_example\_ratio}).
\end{proof}

\begin{proof}[Proof of Proposition~\ref{prop:pk}]
(i) With weights $p_i$, $\sum_ip_i=1$, values $s_i$ and $m=\sum_ip_is_i$:
$\sum_ip_i(s_i-m)^2=\sum_ip_is_i^2-2m\sum_ip_is_i+m^2\sum_ip_i
=\E[S^2]-m^2$ (\lean{secondMoment\_eq\_variance\_add\_sq}).
(ii) For $\lambda>0$, $\rho<1$ the map $m_2\mapsto\lambda m_2/(2(1-\rho))$
has positive slope (\lean{pkWait\_strictMono\_secondMoment}).
(iii) By (i), equal means and $\mathrm{Var}[S_A]<\mathrm{Var}[S_B]$ give
$\E[S_A^2]<\E[S_B^2]$; apply (ii) (\lean{pkWait\_lt\_of\_variance\_lt}).
Example~\ref{ex:pk}: $\E[S_B^2]=\tfrac14(3\cdot10^4+3700^2)=3.43\times10^6$
and the ratio of the two PK expressions is $3.43$ after cancelling the
common denominator (\lean{workloadB\_wait\_ratio}).
\end{proof}

\begin{proof}[Proof of Proposition~\ref{prop:cache}]
(i) $\E[S](q)-\E[S](p)=(q-p)(S^{\mathrm{hit}}-S^{\mathrm{miss}})\le0$ for
$p\le q$; the same identity with squares uses
$0\le S^{\mathrm{hit}}\le S^{\mathrm{miss}}\Rightarrow
(S^{\mathrm{hit}})^2\le(S^{\mathrm{miss}})^2$; and $\rho=\lambda\E[S]$ with
$\lambda\ge0$ (\lean{meanService\_antitone},
\lean{secondMomentService\_antitone}, \lean{utilization\_antitone}).
(ii) Write $\E[W_q](p)=\lambda\E[S^2](p)\,/\,2(1-\rho(p))$. For $p\le q$
with $q\in[0,1]$ and $\rho(p)<1$: the numerator at $q$ is non-negative
(both mixture terms are) and no larger than at $p$ by (i); the denominator
at $q$ is positive and no smaller than at $p$, because $\rho(q)\le\rho(p)<1$.
Hence the ratio at $q$ is at most the ratio at $p$
(\lean{pkWait\_mixture\_antitone}).
Example~\ref{ex:cv2}: for exponential service $\E[S^2]=2\E[S]^2$, so the
ratio of PK expressions is $\E[S^2]/(2\E[S]^2)=(1+\mathrm{CV}^2)/2$ after
cancelling $\lambda/(2(1-\rho))$ (\lean{pkWait\_ratio\_to\_exponential});
the two $\mathrm{CV}^2$ bounds are rational arithmetic
(\lean{mixtureCV2\_agentic\_example}, \lean{mixtureCV2\_cheap\_miss\_example}).
\end{proof}

\begin{proof}[Proof of Proposition~\ref{prop:option}]
Every $a\in A_0$ lies in $A_1$, so $\min_{A_1}J\le J(a)$ for each such $a$,
hence $\min_{A_1}J\le\min_{A_0}J$ (\lean{optimal\_cost\_antitone\_in\_actions};
the memory-controller instance is \lean{enabling\_offload\_never\_hurts}).
For the second claim take $J(\text{keep})=1$, $J(\text{recompute})=5$,
$J(\text{offload})=10$: $\min_{A_0}J\le1<10=J(\text{offload})$
(\lean{always\_offload\_can\_be\_worse}).
\end{proof}

\begin{proof}[Proof of Proposition~\ref{prop:pd}]
(i) Let $m=\min(N_P/s_P,N_D/s_D)$. Then $ms_P\le N_P$ and $ms_D\le N_D$;
adding, $m(s_P+s_D)\le N$, i.e.\ $m\le N/(s_P+s_D)$ (\lean{pd\_le\_agg}).
(ii) Substituting $N_P=Ns_P/(s_P+s_D)$, $N_D=Ns_D/(s_P+s_D)$ makes both
arguments of the $\min$ equal to $N/(s_P+s_D)$
(\lean{pd\_eq\_agg\_at\_rate\_match}); optimality follows from (i)
(\lean{rate\_match\_optimal}).
(iii) $x<\min(a,b,c,d)\iff x<a\wedge x<b\wedge x<c\wedge x<d$
(\lean{pd\_beats\_agg\_iff}); with $g_P=g_D=1$, $I=0$ the compute term is
$N/(s_P+s_D)$ (\lean{pdCompute\_eq\_agg\_of\_no\_gain}).
Example~\ref{ex:pd}: $\mu_A=32/2.5=12.8$; the compute term is
$32/1.5\approx21.3$; so PD wins iff $B_{\mathrm{net}}/\E[K]>12.8$, which
holds at $1000$ and fails at $10$ (\lean{pd\_wins\_example},
\lean{pd\_loses\_example}).
\end{proof}

\begin{proof}[Proof of Proposition~\ref{prop:evict}]
(i) On the sorted list $[4,5,6]$ with $\Delta C=6$ the greedy takes $4$
(remaining $2$) then $5$ and stops: cost $16+25=41$. The set $\{6\}$ frees
$6\ge\Delta C$ at cost $36<41$ (\lean{shortestFirst\_not\_optimal}). The
universal claim, instantiated at this input, would give $41\le36$
(\lean{shortestFirst\_optimality\_claim\_false}). Both are decided by
computation.
(ii) $0.01\cdot10^2=1<4=1\cdot2^2$ (\lean{shortestFirst\_wrong\_with\_resume\_prob}).
\end{proof}

\begin{proof}[Proof of Proposition~\ref{prop:routing}]
(i) $W_1+S_1+F_1<W_2+S_2+M_2+F_2$ rearranges to the stated inequality
(\lean{lookahead\_prefers\_iff}).
(ii) Apply Proposition~\ref{prop:mm1}(iii) with bound $M'=1/\mu+M+F$; the
resulting load $\lambda<\mu$ has $W(\lambda)>W(0)+M+F$
(\lean{affinity\_not\_always\_optimal}).
(iii) $W_D+S_D+I<W_P+S_P+T_{\mathrm{KV}}+W_D$ iff
$S_D+I<W_P+S_P+T_{\mathrm{KV}}$ (\lean{append\_prefill\_rule}).
Example~\ref{ex:routing} is arithmetic (\lean{affinity\_example},
\lean{append\_prefill\_example}).
\end{proof}

\end{document}
